% ============================================================
% Chapter 5 — Conclusion and Future Directions
% ============================================================
\chapter{CONCLUSION AND FUTURE DIRECTIONS}
\label{ch:conclusion}

% ────────────────────────────────────────────
\section{Chapter Overview}

This chapter synthesises the principal findings of this thesis.
Section~\ref{sec:achievements} revisits the five technical objectives stated in Chapter~1 and evaluates the degree to which each was achieved.
Section~\ref{sec:limitations} reflects critically on the methodological and validation limitations that constrain the conclusions that can be drawn from the experimental results.
Section~\ref{sec:future} offers technically precise and actionable recommendations for extending this work over the short and long term.
Section~\ref{sec:closing} provides a final concluding statement.

% ────────────────────────────────────────────
\section{Summary of Achievements Against Objectives}
\label{sec:achievements}

Chapter~1 defined five objectives for this research.
Table~\ref{tab:objectives} maps each objective to its outcome.

\begin{table}[h]
\centering
\caption{Project objectives and their outcomes.}
\label{tab:objectives}
\small
\begin{tabular}{p{0.5\linewidth} p{0.45\linewidth}}
\toprule
\textbf{Objective} & \textbf{Outcome} \\
\midrule
1.\ \textbf{Design} a temporally-extended upsampling architecture incorporating 3D convolutions in the encoder and 3D cross-attention with temporal masking. &
\textit{Achieved.} Both components were designed and integrated into the AnyUp pipeline, producing AnyUp3D. \\
\addlinespace
2.\ \textbf{Construct} the architecture in PyTorch with a modular, backbone-agnostic design compatible with foundation models such as DINOv2 and SigLIP. &
\textit{Achieved with modification.} A fully modular PyTorch implementation was produced. VideoMAE was used as the backbone in lieu of the originally planned image-based foundation models, as its native spatiotemporal tokenization is better suited to the temporal architecture. \\
\addlinespace
3.\ \textbf{Train} the model on the UCF-101 dataset to cultivate temporal priors. &
\textit{Partially achieved.} Training on UCF-101 was completed for 6{,}000 steps. The T-curriculum reached the $T\!=\!4$ stage; the planned $T\!=\!8$ stage was not reached due to compute constraints, and training did not converge to the level achieved by the original 2D AnyUp (100{,}000 steps). \\
\addlinespace
4.\ \textbf{Validate} system performance on the DAVIS-2017 benchmark using J\&F scores and temporal consistency metrics relative to frame-wise baselines. &
\textit{Achieved.} AnyUp3D improved mean temporal cosine similarity from $0.863$ to $0.944$ ($+0.081$) over the 2D AnyUp baseline across all 30 DAVIS clips, and improved mean J\&F from $0.251$ to $0.362$ ($+0.110$) over the bilinear baseline. \\
\addlinespace
5.\ \textbf{Optimize} the inference pipeline to maintain the universal efficiency of the original model while accounting for increased spatio-temporal complexity. &
\textit{Partially achieved.} The T-curriculum strategy was designed to scale complexity incrementally, and mixed-precision inference was employed. Formal inference-time benchmarking (latency and throughput) was not conducted. \\
\bottomrule
\end{tabular}
\end{table}

In summary, the two primary technical contributions — the architectural design and the empirical validation — were fully achieved.
The training objective was partially achieved due to GPU resource constraints, and the optimization objective was addressed at the design level but not formally benchmarked.

% ────────────────────────────────────────────
\section{Limitations and Critical Reflection}
\label{sec:limitations}

\subsection{Methodological Limitations}

\paragraph{Backbone mismatch with dense prediction.}
VideoMAE was chosen because its tubelet embedding natively produces spatiotemporal tokens, making it a principled input for AnyUp3D.
However, VideoMAE is pre-trained for video-level classification, not for dense per-pixel prediction.
Its features consequently lack the spatial specificity that backbones such as Video Swin Transformer or DINOv2 would provide for segmentation tasks.
This limits the absolute level of downstream performance that can be achieved and means the J\&F scores reported in Chapter~4 should be treated as a lower bound on AnyUp3D's potential with a better-aligned backbone.

\paragraph{Heuristic temporal gating.}
The RGB-difference gate in the temporal consistency loss excludes frame pairs whose pixel-level difference is below a fixed threshold.
This heuristic was designed to avoid penalising the model for genuinely different content, but in practice it under-penalises fast-motion sequences such as \textit{bmx-trees} and \textit{motocross-jump}, where large inter-frame RGB differences cause most pairs to be excluded.
The gate thus provides weaker training signal precisely where temporal modelling is hardest, which explains the remaining inconsistency on high-motion clips.

\paragraph{Incomplete T-curriculum.}
The training curriculum was designed in three stages ($T\!=\!2$, $T\!=\!4$, $T\!=\!8$).
Only the first two stages were completed in the available training budget.
Consequently, the model was never exposed to the full eight-frame temporal windows it is architecturally capable of processing, and temporal priors for long-range dependencies were not learned.

\paragraph{Single training run.}
All results are reported from a single training run with a fixed random seed.
While the temporal consistency improvements are observed consistently across all 30 clips — which reduces the likelihood of a chance result — the J\&F scores are not supported by confidence intervals, and the failure cases (\textit{car-shadow}, \textit{bmx-trees}, \textit{dog}) could in principle be seed-dependent.

\subsection{Validation Limitations}

\paragraph{Absence of a spatial-only benchmark.}
No experiment was conducted on a standard image segmentation dataset (e.g., ADE20K) to measure spatial upsampling quality in isolation.
The PCA visualisations in Experiment~1 provide qualitative evidence of spatial structure, and the improvement in $\mathcal{F}$-score in Experiment~3 provides indirect quantitative evidence, but neither constitutes a rigorous, standalone spatial benchmark.
The relative contribution of spatial quality versus temporal consistency to the overall J\&F improvement cannot be quantified from the current results alone.

\paragraph{No ablation study.}
The contribution of each individual design decision — 3D convolutions in the encoder, 3D cross-attention, temporal masking, the RGB-difference gate, the T-curriculum — was not isolated experimentally.
It is therefore not possible to determine which components are responsible for the observed improvements, nor whether any component is redundant.

\paragraph{Single evaluation benchmark.}
The model was validated exclusively on DAVIS-2017.
This is a widely-used standard, but it covers only 30 clips and is biased towards relatively short sequences with a limited range of motion types.
Generalisation to other video benchmarks (e.g., YouTube-VOS, VSPW, or optical flow datasets) has not been demonstrated.

% ────────────────────────────────────────────
\section{Recommendations for Future Research}
\label{sec:future}

\subsection{Short-Term Extensions (1--2 Years)}

\paragraph{Complete the T-curriculum with a larger training budget.}
Extending training to at least 30{,}000 steps — reaching the $T\!=\!8$ stage — would allow the model to learn long-range temporal dependencies and would likely close a substantial portion of the performance gap between AnyUp3D and the original AnyUp.
Training on a single A100 GPU for approximately 48 hours would be sufficient to reach this target.

\paragraph{Conduct a systematic ablation study.}
Isolating the contribution of each component (3D encoder convolutions, 3D cross-attention, temporal masking, RGB-difference gate) through controlled ablation experiments would clarify the relative importance of each design decision.
This would both strengthen the scientific contribution of the work and guide future architectural improvements.

\paragraph{Replace VideoMAE with a dense-prediction backbone.}
Substituting VideoMAE with Video Swin Transformer~\cite{liu2022video} or a frame-level DINOv2 backbone would provide features better aligned to pixel-level tasks.
Based on the J\&F gap between VideoMAE-based pipelines and dedicated VOS methods, a backbone switch alone is estimated to raise the mean J\&F by $0.10$--$0.15$ without any changes to AnyUp3D itself.

\paragraph{Validate on additional benchmarks.}
Evaluating AnyUp3D on YouTube-VOS~\cite{xu2018youtube} and a standard image segmentation dataset such as ADE20K would demonstrate generality beyond DAVIS-2017 and provide the spatial-only baseline that was absent from the current evaluation.

\paragraph{Replace the heuristic RGB-difference gate with a learned temporal weighting.}
Training a lightweight network to predict the reliability of each frame pair as a training signal — rather than using a fixed RGB-difference threshold — would provide more principled temporal weighting and address the under-penalisation of fast-motion clips identified in Section~\ref{sec:limitations}.

\subsection{Long-Term Research Avenues (3--5 Years)}

\paragraph{Extend universality to additional modalities.}
The backbone-agnostic design of AnyUp3D is not limited to RGB video.
Extending the framework to handle features from depth sensors, LiDAR point clouds, or thermal cameras would broaden its applicability to autonomous driving and robotics perception pipelines, where multi-modal temporal consistency is a persistent open challenge.

\paragraph{Efficient temporal attention for real-time deployment.}
The current 3D cross-attention mechanism has quadratic complexity in the temporal dimension.
Replacing it with linear attention or sparse windowed attention — analogous to the approach taken in Video Swin Transformer — would reduce inference time sufficiently to enable real-time processing at 30 frames per second on consumer-grade GPUs.
This is a prerequisite for deployment in video editing tools and live streaming applications.

\paragraph{Integration with video diffusion models.}
Large-scale video generation models such as Sora and Stable Video Diffusion produce temporally coherent latent representations that are structurally similar to the feature maps AnyUp3D is designed to upsample.
Integrating AnyUp3D as a temporal upsampling stage within a video diffusion pipeline could improve the spatial resolution of generated videos without retraining the generative backbone, representing a computationally efficient path to higher-fidelity video synthesis.

\paragraph{Theoretical analysis of temporal stability guarantees.}
The current work demonstrates empirical temporal consistency improvements but does not provide a theoretical characterisation of the conditions under which the 3D cross-attention mechanism is guaranteed to produce stable features.
Deriving bounds on the inter-frame feature variation as a function of architecture hyperparameters (window size, number of heads, temporal kernel size) would lay a rigorous foundation for deploying AnyUp3D in safety-critical applications such as medical video analysis.

% ────────────────────────────────────────────
\section{Final Concluding Remarks}
\label{sec:closing}

This thesis demonstrated that the resolution gap inherent to Vision Foundation Models can be addressed in the temporal domain as well as the spatial domain.
By introducing 3D convolutional encoding and 3D cross-attention with temporal masking into the AnyUp framework, AnyUp3D produces high-resolution feature maps that are stable across time — achieving a mean temporal cosine similarity improvement of $+0.081$ over frame-independent upsampling on every clip in the DAVIS-2017 benchmark, and translating this into a $+0.110$ improvement in downstream video object segmentation quality.
The work establishes a principled, backbone-agnostic foundation for temporally consistent universal upsampling, confirming that explicitly modelling inter-frame dependencies is both necessary and sufficient to overcome the flickering artifacts that limit the practical utility of image-based upsampling methods in video applications.
